\chapter{Time and Space}\label{ch:timespace}

\begin{center}\itshape
With more time you can solve more problems, and with more memory too;\\
but guessing helps memory much less than anyone expected.\\
\upshape(the content of this chapter, in our words)
\end{center}

Computability theory sorts problems into two kinds: those a machine can solve, and the others. Its answers do not care whether the machine needs a second or the age of the universe. In 1965 two young researchers at the General Electric Research Laboratory in Schenectady, Juris Hartmanis and Richard Stearns, published a paper with the title \emph{On the computational complexity of algorithms}. They measured the number of steps a Turing machine needs as a function of the length of its input, defined the class of problems solvable within a given time bound, and proved that more time solves strictly more problems. The paper gave its name to a field, and in 1993 its authors received the Turing Award for it.

The same year, Alan Cobham and Jack Edmonds argued, independently, that ``solvable in a number of steps bounded by a polynomial in the length of the input'' is the right mathematical notion of an efficient algorithm. Edmonds, whose blossom algorithm for matchings is in Tome~C2, Chapter~6, called such algorithms ``good''. A polynomial bound is robust: it does not depend on the machine model, since the models simulate each other with polynomial overhead (Theorem~\ref{thm:multitape}), and polynomials are closed under composition, so a good algorithm calling a good subroutine is good.

Memory turned out to behave differently from time. In 1970 Walter Savitch showed that a nondeterministic machine, which may guess, can be simulated deterministically with only the square of its memory, whereas the best known simulation for time is exponential. Then, in 1987, came a surprise. In 1964 S.-Y. Kuroda had asked whether the complement of a context-sensitive language is context-sensitive, a question about nondeterministic linear memory that resisted for twenty-three years, while nearly everyone expected the answer to be no. Neil Immerman, in the United States, and Róbert Szelepcsényi, a student of twenty-one in Bratislava, proved independently, with essentially the same idea, that the answer is yes. They shared the Gödel Prize in 1995.

\begin{think}
To decide whether one vertex of a graph with a million vertices can reach another, how much memory do you need, if you may not modify the graph? Is it easier to convince someone that a path exists, or that no path exists?
\end{think}

\section{Measuring resources}\label{sec:resources}

We use the multitape Turing machines of Chapter~\ref{ch:turing}. A machine \emph{decides} a language if it halts on every input, accepting exactly the words of the language. It \emph{runs in time} $t(n)$ if on every input of length $n$ it halts after at most $t(n)$ steps; a nondeterministic machine runs in time $t(n)$ if all its computations halt within $t(n)$ steps, and accepts a word if one of them accepts. For memory, a machine has a read-only \emph{input tape}, on which the input is written between two end markers, and work tapes; it \emph{runs in space} $s(n)$ if on inputs of length $n$ it visits at most $s(n)$ cells of its work tapes, along every computation. The input does not count, so that space below $n$ makes sense.

\begin{definition}\label{def:classes}
For $f:\N\to\N$, $\mathsf{DTIME}(f)$ is the class of languages decided by deterministic multitape machines in time $O(f(n))$, and $\mathsf{NTIME}(f)$, $\mathsf{DSPACE}(f)$, $\mathsf{NSPACE}(f)$ are defined in the same way. Then
\[
\mathsf L=\mathsf{DSPACE}(\log n),\quad\mathsf{NL}=\mathsf{NSPACE}(\log n),\quad\mathsf P=\bigcup_k\mathsf{DTIME}(n^k),\quad\mathsf{NP}=\bigcup_k\mathsf{NTIME}(n^k),
\]
\[
\mathsf{PSPACE}=\bigcup_k\mathsf{DSPACE}(n^k),\quad\mathsf{NPSPACE}=\bigcup_k\mathsf{NSPACE}(n^k),\quad\mathsf{EXP}=\bigcup_k\mathsf{DTIME}(2^{n^k}).
\]
For a class $\mathsf C$, $\mathsf{coC}$ is the class of complements of languages in $\mathsf C$.
\end{definition}

Logarithmic space is the space of a constant number of pointers into the input, or counters up to the input length: with it one can check that a word is a palindrome, compare two numbers, or follow a path in a graph given as input. A graph problem in $\mathsf L$ can be solved by a program which may read the graph but keeps only a few vertex names in memory.

The main tool for space is the \emph{configuration}: on an input $x$ of length $n$, a configuration of a machine with $k$ work tapes using space $s$ records the state, the position of the input head ($n+2$ choices), the positions of the work heads and the contents of the work tapes. There are at most
\begin{equation}\label{eq:configs}
|Q|\cdot(n+2)\cdot s^k\cdot|\Gamma|^{ks}\le2^{c(s+\log n)}
\end{equation}
of them, for a constant $c$ depending on the machine. The \emph{configuration graph} of the machine on $x$ has the configurations as vertices and an arrow from each to its possible successors; the machine accepts $x$ if and only if some accepting configuration is reachable from the initial one. Whether $C\to C'$ is an arrow can be checked in space $O(s+\log n)$, by looking at the two configurations and at the input symbol under the head.

\begin{proposition}\label{prop:inclusions}
For $f(n)\ge n$, $\mathsf{NTIME}(f)\subseteq\mathsf{DSPACE}(f)$ when $f$ is computable in space $O(f)$; and for $s(n)\ge\log n$, $\mathsf{NSPACE}(s)\subseteq\bigcup_c\mathsf{DTIME}(2^{cs})$. Consequently
\[
\mathsf L\subseteq\mathsf{NL}\subseteq\mathsf P\subseteq\mathsf{NP}\subseteq\mathsf{PSPACE}\subseteq\mathsf{NPSPACE}\subseteq\mathsf{EXP}.
\]
\end{proposition}

\begin{proof}
A nondeterministic machine running in time $O(f)$ makes at most $O(f(n))$ choices, each among a bounded number $d$ of moves. A deterministic machine can try all sequences of choices, in the order of their numbering, keeping on a tape the current sequence ($O(f(n))$ cells, after computing $f(n)$) and simulating the computation it describes, which visits $O(f(n))$ cells; it reuses the same cells for each sequence, and accepts if one of them accepts.

For the second inclusion, the configuration graph has at most $2^{c(s(n)+\log n)}\le2^{2cs(n)}$ vertices. A deterministic machine can list them and search the graph (breadth-first, Tome~C2, Chapter~4) in time polynomial in its size, which is $2^{O(s(n))}$. The chain follows: deterministic machines are special nondeterministic ones, $\mathsf{NL}\subseteq\mathsf{DTIME}(2^{O(\log n)})=\mathsf P$, $\mathsf{NP}\subseteq\mathsf{PSPACE}$ by the first part, and $\mathsf{NPSPACE}\subseteq\mathsf{DTIME}(2^{O(n^k)})\subseteq\mathsf{EXP}$.
\end{proof}

Which of these inclusions are strict? Almost nothing is known. Section~\ref{sec:hierarchies} shows that $\mathsf P\ne\mathsf{EXP}$ and $\mathsf{NL}\ne\mathsf{PSPACE}$, so at least one inclusion between $\mathsf P$ and $\mathsf{EXP}$ is strict, and at least one between $\mathsf{NL}$ and $\mathsf{PSPACE}$; Section~\ref{sec:savitch} shows that $\mathsf{PSPACE}=\mathsf{NPSPACE}$. Nobody knows which of the others are equalities; Figure~\ref{fig:classes} shows the state of knowledge.

\begin{figure}[ht]
\centering
\fig{ch08-classes}
\caption{The classes of this chapter, as they are believed to be: every inclusion strict, and $\mathsf{NP}\ne\mathsf{coNP}$. Only $\mathsf P\ne\mathsf{EXP}$ and $\mathsf{NL}\ne\mathsf{PSPACE}$ (and their consequences) are proved; for all we know, $\mathsf L=\mathsf{NP}$, or $\mathsf P=\mathsf{PSPACE}$.}
\label{fig:classes}
\end{figure}

\section{Hierarchies}\label{sec:hierarchies}

To show that more time gives more power, we let a machine with more time disagree with every machine with less time, by Cantor's diagonal argument (Theorem~\ref{thm:halting}). The diagonal machine must know when to stop, so the bounds must themselves be computable within the resources they allow.

\begin{definition}\label{def:constructible}
A function $g:\N\to\N$ with $g(n)\ge n$ is \emph{time-constructible} if some machine, on input $1^n$, writes $g(n)$ in binary in time $O(g(n))$. A function $g$ with $g(n)\ge\log n$ is \emph{space-constructible} if some machine, on input $1^n$, writes $g(n)$ in binary in space $O(g(n))$.
\end{definition}

The functions $n^k$, $n\log n$, $2^n$, $\log n$ and all the functions one meets are constructible (Exercise~\ref{ex:constructible}). The diagonal machine also needs to simulate other machines efficiently.

\begin{lemma}\label{lem:efficientuniversal}
There is a multitape machine $U$ which, on input $\langle M\rangle$ and $w$, simulates the multitape machine $M$ on $w$, performing $O(c_Mt^2)$ steps for $t$ steps of $M$, where $c_M$ depends only on $M$. A variant uses space $O(c_M's)$ when $M$ uses space $s$, and $O(\log n)$ for the position of the input head.
\end{lemma}

\begin{proof}
By Theorem~\ref{thm:multitape}, $M$ is simulated by a one-tape machine $S$, computable from $\langle M\rangle$, with $O(t^2)$ steps for $t$ steps of $M$. The machine $U$ writes $\langle S\rangle$ on a work tape, keeps the tape of $S$ on a second work tape and the current state of $S$ on a third, and simulates a step of $S$ by looking up the matching transition in $\langle S\rangle$, in $O(|\langle S\rangle|)$ steps. For space, $U$ keeps instead the $k$ work tapes of $M$ on one work tape, each symbol of $M$ coded by a block of $\lceil\log_2|\Gamma_M|\rceil$ bits, and reads $M$'s input on its own input tape, keeping the position of $M$'s input head in binary.
\end{proof}

\begin{theorem}[time hierarchy; Hartmanis and Stearns, 1965]\label{thm:timehierarchy}
If $g$ is time-constructible and $f(n)^2=o(g(n))$, then $\mathsf{DTIME}(f)\subsetneq\mathsf{DTIME}(g)$.
\end{theorem}

\begin{proof}
Consider the machine $D$ which, on an input $w$ of length $n$:
\begin{enumerate}
\item computes $g(n)$ in binary on a separate \emph{clock} tape (time $O(g(n))$), and rejects if $w$ is not of the form $\langle M\rangle10^k$ for a machine $M$;
\item simulates $M$ on $w$ with the machine $U$ of Lemma~\ref{lem:efficientuniversal}, decreasing the clock by one at each step of $U$, and rejects if the clock reaches $0$;
\item if the simulation finishes in time, accepts if $M$ rejects $w$, and rejects if $M$ accepts $w$.
\end{enumerate}
Decreasing a binary counter costs $O(1)$ steps on average (Tome~C2, Chapter~3), so $D$ runs in time $O(g(n))$ and $L(D)\in\mathsf{DTIME}(g)$. Suppose $L(D)$ were decided by a machine $M$ running in time $af(n)$. Its simulation on an input $w=\langle M\rangle10^k$ takes at most $bc_Ma^2f(n)^2$ steps of $U$, for a constant $b$. As $f(n)^2=o(g(n))$, this is less than $g(n)$ when $n$ is large, which we can ensure by taking $k$ large. Then $D$ completes the simulation and gives the opposite answer to $M$ on $w$, so $L(D)\ne L(M)$.
\end{proof}

With a better simulation, by Fred Hennie and Stearns in 1966, two tapes simulate any number of tapes with a logarithmic overhead, and the hypothesis becomes $f(n)\log f(n)=o(g(n))$ (quoted). Ours is enough for the main corollary: $\mathsf{DTIME}(n^k)\subsetneq\mathsf{DTIME}(n^{2k+1})$ for every $k$, so there is no largest polynomial; and since $\mathsf P\subseteq\mathsf{DTIME}(2^n)$ fails to contain the diagonal language for $g(n)=2^{2n}$, while every $n^k$ satisfies $(n^k)^2=o(2^n)$,
\[
\mathsf P\subsetneq\mathsf{DTIME}(2^n)\subsetneq\mathsf{EXP}.
\]

\begin{theorem}[space hierarchy; Stearns, Hartmanis and Lewis, 1965]\label{thm:spacehierarchy}
If $g$ is space-constructible and $f(n)=o(g(n))$, then $\mathsf{DSPACE}(f)\subsetneq\mathsf{DSPACE}(g)$.
\end{theorem}

\begin{proof}
The machine $D$, on an input $w$ of length $n$, computes $g(n)$ and marks $g(n)$ cells of a work tape; it rejects if $w$ is not of the form $\langle M\rangle10^k$. It simulates $M$ on $w$ as in Lemma~\ref{lem:efficientuniversal}, within the marked cells, rejecting if the simulation needs more. It also counts the simulated steps in binary on $2g(n)+2$ cells, and accepts if the counter overflows. If $M$ halts within these limits, $D$ gives the opposite answer. Then $D$ runs in space $O(g(n))$.

Suppose $M$ decides $L(D)$ in space $af(n)$. Its simulation needs $c_M'af(n)$ cells, fewer than $g(n)$ for large $n$. By \eqref{eq:configs}, $M$ has at most $2^{c(af(n)+\log n)}$ configurations on inputs of length $n$; as $f=o(g)$ and $\log n\le g(n)$, this is less than $2^{2g(n)+2}$ for large $n$. A deterministic computation that runs longer than the number of configurations repeats one, and loops forever; $M$ always halts, so it halts before the counter overflows. Taking $k$ large, $D$ completes the simulation on $\langle M\rangle10^k$ and disagrees with $M$.
\end{proof}

So $\mathsf L\subsetneq\mathsf{PSPACE}$, and $\mathsf{DSPACE}(n^k)\subsetneq\mathsf{DSPACE}(n^{k+1})$. The same diagonal argument works for nondeterministic time, but it is much harder, since a nondeterministic machine cannot simply flip the answer of another one: complementing seems to need exponential time. The nondeterministic time hierarchy was proved by Stephen Cook in 1972 and in its sharp form by Joel Seiferas, Michael Fischer and Albert Meyer in 1978, with a clever delayed diagonalisation (quoted).

\section{Guessing with little memory}\label{sec:savitch}

The basic problem of nondeterministic space is reachability: $\mathrm{PATH}$ is the language of directed graphs $G$ with two vertices $s,t$ such that $t$ is reachable from $s$. It is in $\mathsf{NL}$: guess a path one vertex at a time, keeping only the current vertex and a step counter, and stop after $n$ steps. And every problem in $\mathsf{NSPACE}(s)$ is a reachability problem in the configuration graph, a graph with $2^{O(s)}$ vertices which is never written down: its arrows are computed on demand.

\begin{theorem}[Savitch, 1970]\label{thm:savitch}
For space-constructible $s(n)\ge\log n$, $\mathsf{NSPACE}(s)\subseteq\mathsf{DSPACE}(s^2)$. In particular $\mathsf{PSPACE}=\mathsf{NPSPACE}$ and $\mathsf{NL}\subseteq\mathsf{DSPACE}(\log^2n)$.
\end{theorem}

\begin{proof}
Let $N$ be a nondeterministic machine running in space $s$. We may assume that $N$ has a single accepting configuration $C_{\mathrm{acc}}$ on each input, by making it erase its work tapes and move its heads to the left end before accepting. By \eqref{eq:configs} there are at most $2^{ds(n)}$ configurations, for a constant $d$, so $N$ accepts $x$ if and only if $C_{\mathrm{acc}}$ can be reached from the initial configuration $C_0$ in at most $2^{ds(n)}$ steps. Consider the recursive procedure:
\begin{quote}
$\mathrm{Reach}(C,C',t)$: if $t=1$, return whether $C=C'$ or $C\to C'$ is an arrow. Otherwise, for every configuration $C''$: if $\mathrm{Reach}(C,C'',t/2)$ and $\mathrm{Reach}(C'',C',t/2)$, return true. Return false.
\end{quote}
For $t$ a power of $2$, it returns whether $C'$ can be reached from $C$ in at most $t$ steps: a path of length at most $t$ has a middle configuration $C''$ splitting it into two paths of length at most $t/2$ (Figure~\ref{fig:savitch}). The deterministic machine computes $\mathrm{Reach}(C_0,C_{\mathrm{acc}},2^{ds(n)})$. The recursion has depth $ds(n)$, and each level of the stack stores two configurations, a candidate $C''$ and the exponent of $t$, in $O(s(n))$ cells; the two recursive calls reuse the same space one after the other. The total is $O(s(n)^2)$.
\end{proof}

\begin{figure}[ht]
\centering
\fig{ch08-savitch}
\caption{Savitch's recursion. To decide whether $c_2$ is reachable from $c_1$ in $t=8$ steps, try every candidate midpoint $m$ and decide, one after the other, the two problems with $t=4$, and so on. The stack holds one frame for each level: $\log_2t$ frames of $O(s)$ cells each. The time is enormous (each level tries all configurations), but the space is small.}
\label{fig:savitch}
\end{figure}

For time, the analogous question, whether $\mathsf P=\mathsf{NP}$, is the subject of the next chapter. For space the price of nondeterminism is at most a square, and the classes $\mathsf{PSPACE}$ and $\mathsf{NPSPACE}$ coincide. Whether $\mathsf L=\mathsf{NL}$ is open. A special case was settled by Omer Reingold in 2005: reachability in \emph{undirected} graphs is in $\mathsf L$ (quoted); his algorithm walks along the graph after transforming it into an expander, the graphs of Tome~C1, Chapter~9.

\begin{theorem}[Immerman, 1988; Szelepcsényi, 1988]\label{thm:immerman}
The complement of $\mathrm{PATH}$ is in $\mathsf{NL}$. Consequently $\mathsf{NL}=\mathsf{coNL}$, and more generally $\mathsf{NSPACE}(s)=\mathsf{coNSPACE}(s)$ for space-constructible $s(n)\ge\log n$.
\end{theorem}

The difficulty is that a nondeterministic machine proves membership by exhibiting something, a path; to prove that $t$ is \emph{not} reachable it must exhibit an absence. The trick is to count. If the machine knows how many vertices are reachable from $s$, it can check that $t$ is not one of them: guess all of them, with paths, check that there are that many, and that $t$ is not among them. And the number itself can be computed in the same way, layer after layer (Figure~\ref{fig:counting}).

\begin{figure}[ht]
\centering
\fig{ch08-counting}
\caption{Inductive counting. The vertices at distance $0,1,2,3,4$ from $s$ are shaded from dark to light; $R_i$ is the set of vertices at distance at most $i$, and $c_i=|R_i|$. Knowing $c_i$, a nondeterministic machine can list $R_i$ with certificates, and so decide membership in $R_{i+1}$ in both directions. The count stabilises at $c_4=7$, and $t$, which has arrows out but none in, is not among the seven.}
\label{fig:counting}
\end{figure}

\begin{proof}
Let $G$ have vertices $1,\dots,n$, and let $R_i$ be the set of vertices reachable from $s$ by a path of length at most $i$, and $c_i=|R_i|$. We describe a nondeterministic procedure using $O(\log n)$ space, in which some computations \emph{give up} (reject), and every computation that does not give up produces the correct result. We use three subroutines.

\emph{Guess a path.} Given $u$ and $i$, guess a path from $s$ of length at most $i$, vertex by vertex, and give up if it does not end at $u$. This uses a counter and the current vertex.

\emph{Enumerate $R_i$, given $c_i$.} For each vertex $u=1,\dots,n$ in turn, guess whether $u\in R_i$; if the guess is yes, confirm it by guessing a path to $u$ of length at most $i$, and count it. At the end, give up unless the count is $c_i$. A computation that does not give up has confirmed $c_i$ members of $R_i$, hence all of them, in increasing order, and it has listed each $u\in R_i$ and no other vertex.

\emph{Compute $c_{i+1}$ from $c_i$.} For each vertex $v$, decide whether $v\in R_{i+1}$: enumerate $R_i$ as above, and answer yes if some listed $u$ equals $v$ or has an arrow $u\to v$; this is correct, since $R_{i+1}$ consists of $R_i$ and its successors. Count the $v$ with answer yes.

The machine starts from $c_0=1$, computes $c_1,\dots,c_{n-1}$, keeping only the last one, and $R_{n-1}$ is the set of all vertices reachable from $s$. Finally it enumerates $R_{n-1}$ and accepts if $t$ is not listed. Some computation guesses everything correctly and accepts if $t$ is unreachable; no computation accepts if $t$ is reachable, since an enumeration which does not give up lists $t$. The memory holds a constant number of vertices and counters: $O(\log n)$.

For a language in $\mathsf{NSPACE}(s)$, apply the same procedure to the configuration graph on the input $x$, with the initial configuration as $s$ and the accepting configuration as $t$, computing arrows on demand: the vertices are now names of $O(s)$ cells, the counters go up to $2^{O(s)}$, and the procedure accepts exactly the $x$ rejected by the original machine.
\end{proof}

Context-sensitive grammars, which allow rules $uAv\to uwv$ that rewrite a symbol in a context, generate exactly the languages in $\mathsf{NSPACE}(n)$, by a theorem of Kuroda (1964, quoted). So Kuroda's question, open for twenty-three years, was answered: the complement of a context-sensitive language is context-sensitive.

\section{Oracles and the limits of diagonalisation}\label{sec:oracles}

Every separation in this chapter was proved by simulation and diagonalisation. Could the same methods separate $\mathsf P$ from $\mathsf{NP}$? In 1975 Theodore Baker, John Gill and Robert Solovay showed that they cannot, by giving the machines an oracle.

An \emph{oracle machine} with oracle $A\subseteq\{0,1\}^*$ has an extra tape on which it can write a word $y$, then enter a query state, and in one step learn whether $y\in A$. The classes $\mathsf P^A$ and $\mathsf{NP}^A$ are defined like $\mathsf P$ and $\mathsf{NP}$ with oracle machines. The proofs of this chapter \emph{relativise}: they work word for word with oracles, since simulating an oracle machine only needs the same oracle. A proof that $\mathsf P\ne\mathsf{NP}$ by such methods would therefore prove $\mathsf P^A\ne\mathsf{NP}^A$ for every oracle $A$, and a proof of $\mathsf P=\mathsf{NP}$ would prove $\mathsf P^A=\mathsf{NP}^A$ for every $A$.

\begin{theorem}[Baker, Gill and Solovay, 1975]\label{thm:bgs}
There are oracles $A$ and $B$ with $\mathsf P^A=\mathsf{NP}^A$ and $\mathsf P^B\ne\mathsf{NP}^B$.
\end{theorem}

\begin{proof}
\emph{The oracle $A$.} Let $A$ be the set of words $\langle M,x,1^k\rangle$ such that the deterministic machine $M$ accepts $x$ using at most $k$ cells of work space. Every $L\in\mathsf{PSPACE}$, decided by $M$ in space $n^c$ say, is in $\mathsf P^A$: on input $x$, ask whether $\langle M,x,1^{|x|^c}\rangle\in A$. Conversely, $\mathsf{NP}^A\subseteq\mathsf{PSPACE}$: a nondeterministic oracle machine running in time $n^c$ asks queries of length at most $n^c$; a deterministic machine can try all its sequences of choices one after the other in polynomial space, and answer each query $\langle M,x,1^k\rangle$ itself by simulating $M$ on $x$ in space $k$, with a step counter up to the number of configurations \eqref{eq:configs} to detect loops, in space polynomial in the length of the query. So $\mathsf{NP}^A\subseteq\mathsf{PSPACE}\subseteq\mathsf P^A\subseteq\mathsf{NP}^A$.

\emph{The oracle $B$.} For any $B$, let $L_B=\{1^n:B\text{ contains a word of length }n\}$. It is in $\mathsf{NP}^B$: guess a word of length $n$ and ask the oracle. We build $B$ so that $L_B\notin\mathsf P^B$. List the pairs $(M_i,k_i)$ of deterministic oracle machines and integers; every language in $\mathsf P^B$ is decided by some $M_i$ within $n^{k_i}+k_i$ steps, and we stop $M_i$ after that many steps. We decide, in stages, which words belong to $B$; at each stage finitely many words have been decided. At stage $i$, choose $n$ larger than the length of every word decided so far and with $2^n>n^{k_i}+k_i$. Run $M_i$ on $1^n$ for $n^{k_i}+k_i$ steps, answering its queries according to the decisions so far, and ``no'' for undecided words, which we then decide to be outside $B$. If $M_i$ accepts, decide that $B$ has no word of length $n$, so $1^n\notin L_B$. If $M_i$ rejects, it has queried fewer than $2^n$ words; put into $B$ a word of length $n$ it did not query, so $1^n\in L_B$. In both cases the answers given to $M_i$ remain correct for the final $B$, since later stages only decide longer words, so $M_i$ with oracle $B$ errs on $1^n$. No machine decides $L_B$ in polynomial time with oracle $B$.
\end{proof}

So simulation and diagonalisation alone cannot settle $\mathsf P$ versus $\mathsf{NP}$: any proof must use something specific about computation without oracles. Two other barriers were found later, \emph{natural proofs} in 1994 (Chapter~12) and \emph{algebrisation} in 2008, and they block most other known techniques. The problem is not only unsolved; we know a good deal about why it is hard.

\begin{beauty}
Immerman and Szelepcsényi's proof turns a question about absence into a question about counting: to certify that something is not there, count everything that is, and show it. The same idea, counting as a certificate, runs through mathematics, from the proof that a polynomial has no more roots than its degree to the class number formula. Here it overturned a conjecture that almost everyone believed, and did it with a proof that fits on a page.
\end{beauty}

\section{Exercises}

\begin{exercise}[$\star$]\label{ex:logspace}
Show that the palindromes over $\{a,b\}$ and the language $\{a^nb^nc^n:n\ge0\}$ are in $\mathsf L$.
\end{exercise}

\begin{exercise}[$\star\star$]\label{ex:logarith}
Show that the sum and the product of two numbers given in binary can be computed in logarithmic space: there is a machine with a write-only output tape, moving only to the right, which writes the result using $O(\log n)$ cells of work space.
\end{exercise}

\begin{exercise}[$\star$]\label{ex:hierarchyexamples}
Show that $\mathsf{DTIME}(n)\subsetneq\mathsf{DTIME}(n^3)$, $\mathsf{DSPACE}(\log n)\subsetneq\mathsf{DSPACE}(\log^2n)$ and $\mathsf{NL}\subsetneq\mathsf{PSPACE}$.
\end{exercise}

\begin{exercise}[$\star\star$]\label{ex:constructible}
Show that $n^2$, $2^n$ and $n\lfloor\log_2n\rfloor$ are time-constructible, and $\lfloor\log_2n\rfloor$ and $\lceil\sqrt n\rceil$ space-constructible.
\end{exercise}

\begin{exercise}[$\star\star$]\label{ex:padding}
Let $\mathsf{NEXP}=\bigcup_k\mathsf{NTIME}(2^{n^k})$. Show that if $\mathsf P=\mathsf{NP}$, then $\mathsf{EXP}=\mathsf{NEXP}$. (Pad the inputs: for $L\in\mathsf{NTIME}(2^{n^k})$, consider $\{x\#1^{2^{|x|^k}}:x\in L\}$.)
\end{exercise}

\begin{exercise}[$\star$]\label{ex:dag}
Show that the language of directed graphs without a cycle is in $\mathsf{NL}$.
\end{exercise}

\begin{exercise}[$\star\star$]\label{ex:twosatnl}
Show that 2-SAT, the satisfiable formulas in conjunctive normal form with two literals per clause, is in $\mathsf{NL}$. (Use the implication graph of Tome~C1, Theorem~1.17 and the theorem of Immerman and Szelepcsényi.)
\end{exercise}

\begin{exercise}[$\star\star$]\label{ex:bipartite}
Show that the language of bipartite (undirected) graphs is in $\mathsf{NL}$.
\end{exercise}

\begin{exercise}[$\star\star$]\label{ex:bgsconp}
Adapt the construction of the oracle $B$ to build an oracle $C$ with $\mathsf{NP}^C\ne\mathsf{coNP}^C$. (Use $\{1^n:\text{every word of length }n\text{ is in }C\}$, which is in $\mathsf{coNP}^C$, and diagonalise against nondeterministic oracle machines.)
\end{exercise}

\begin{exercise}[$\star\star\star$]\label{ex:pnotlinspace}
Show that $\mathsf P\ne\mathsf{DSPACE}(n)$, although nobody knows whether either contains the other. (If they were equal, a padding argument would put $\mathsf{DSPACE}(n^2)$ inside $\mathsf P$.)
\end{exercise}

\section{Hints}

\begin{hint}{ex:logspace}
Keep two indices into the input, in binary.
\end{hint}
\begin{hint}{ex:logarith}
Compute the result digit by digit from the right, recomputing the carry when needed; for the product, the $k$-th digit depends on the sums $\sum_{i+j=k}a_ib_j$ and a carry which is at most $n$.
\end{hint}
\begin{hint}{ex:hierarchyexamples}
Apply the hierarchy theorems, and Savitch's theorem for the last one.
\end{hint}
\begin{hint}{ex:constructible}
Use binary counters; for $\sqrt n$, try $m=1,2,\dots$ and compare $m^2$ with $n$ in binary.
\end{hint}
\begin{hint}{ex:padding}
The padded language is in $\mathsf{NP}$, hence in $\mathsf P$; then remove the padding.
\end{hint}
\begin{hint}{ex:dag}
A graph has a cycle if and only if some vertex can reach itself by a path of length between $1$ and $n$. That is in $\mathsf{NL}$; then use Theorem~\ref{thm:immerman}.
\end{hint}
\begin{hint}{ex:twosatnl}
The formula is unsatisfiable if and only if for some variable $x$, $\neg x$ is reachable from $x$ and $x$ from $\neg x$.
\end{hint}
\begin{hint}{ex:bipartite}
A graph is bipartite if and only if it has no cycle of odd length. Guess an odd closed walk.
\end{hint}
\begin{hint}{ex:bgsconp}
Run $N_i$ on $1^n$ with every undecided word of length $n$ in the oracle. If some computation accepts, freeze the fewer than $2^n$ queries it makes and remove one unqueried word of length $n$; if none accepts, keep all words of length $n$.
\end{hint}
\begin{hint}{ex:pnotlinspace}
Show that $\mathsf{DSPACE}(n)=\mathsf P$ would imply $\mathsf{DSPACE}(n^2)\subseteq\mathsf P$ by padding, and contradict the space hierarchy theorem.
\end{hint}

\section{Solutions}

\begin{solution}{ex:logspace}
\idea Logarithmic space is a constant number of pointers into the input.
\writeup For palindromes: for $i=1,\dots,\lfloor n/2\rfloor$, compare the $i$-th and $(n+1-i)$-th letters. The machine keeps $i$ and $n$ in binary ($O(\log n)$ cells) and finds the $j$-th letter by moving the input head to the left marker and then right while counting up to $j$. For $a^nb^nc^n$: check with one pass that the word is in $a^*b^*c^*$ (a finite automaton, no memory), then count the $a$'s, $b$'s and $c$'s in three binary counters and compare.
\end{solution}

\begin{solution}{ex:logarith}
\idea Output digits as they are found, and recompute what you cannot store.
\writeup Let the inputs be $a=\sum_ia_i2^i$ and $b=\sum_ib_i2^i$ with $n$ digits. For the sum, the carry into position $k$ is $0$ or $1$; the machine scans $k=0,1,\dots$, reading $a_k$ and $b_k$ by moving the input head with a binary counter, keeps the carry, and writes $a_k+b_k+\text{carry}\bmod2$. It writes the digits from the least significant one; that is the convention of the output. For the product, the $k$-th digit of $ab$ is $(s_k+r_k)\bmod2$, where $s_k=\sum_{i+j=k}a_ib_j\le n$ and $r_k$ is the carry from position $k-1$, given by $r_{k+1}=\lfloor(s_k+r_k)/2\rfloor$. By induction $r_k\le n$: if $r_k\le n$ then $r_{k+1}\le(n+n)/2=n$. So the machine keeps $k$, $r_k$ and computes $s_k$ with an inner loop over $i$, all numbers at most $2n$, in $O(\log n)$ cells.
\end{solution}

\begin{solution}{ex:hierarchyexamples}
\idea Each separation is one hierarchy theorem, with room to spare.
\writeup $g(n)=n^3$ is time-constructible and $(n)^2=o(n^3)$, so Theorem~\ref{thm:timehierarchy} applies. $\log^2n$ is space-constructible (compute $\lfloor\log_2n\rfloor$ by counting the digits of $n$, then square it) and $\log n=o(\log^2n)$, so Theorem~\ref{thm:spacehierarchy} applies. By Theorem~\ref{thm:savitch}, $\mathsf{NL}\subseteq\mathsf{DSPACE}(\log^2n)\subsetneq\mathsf{DSPACE}(n)\subseteq\mathsf{PSPACE}$.
\end{solution}

\begin{solution}{ex:constructible}
\idea Arithmetic in binary is cheap: everything is a loop over a counter.
\writeup On input $1^n$, a machine computes $n$ in binary by incrementing a counter at each input letter, in $O(n)$ steps (amortised, Tome~C2, Chapter~3). Then $n^2$ in binary by the school multiplication, $O(\log^2n)$ steps; $2^n$ in binary is a $1$ followed by $n$ zeros, $O(n)$ steps; $\lfloor\log_2n\rfloor$ is the number of digits of $n$ minus one, and $n\lfloor\log_2n\rfloor$ a product; all in time $O(n)\le O(g(n))$. For space: $\lfloor\log_2n\rfloor$ is computed in space $O(\log n)$ by the counter. For $\lceil\sqrt n\rceil$, try $m=1,2,\dots$, computing $m^2$ in binary and comparing it with $n$; the numbers have $O(\log n)$ digits, which is $O(\sqrt n)$.
\end{solution}

\begin{solution}{ex:padding}
\idea Padding makes an exponential-time problem look polynomial in a longer input; a collapse among small classes then propagates upwards.
\writeup Let $L\in\mathsf{NTIME}(2^{n^k})$, decided by $N$, and $L'=\{x\#1^{2^{|x|^k}}:x\in L\}$, of length $m=|x|+1+2^{|x|^k}$. A nondeterministic machine decides $L'$ in time $O(m)$: check the format (computing $2^{|x|^k}$ in time $O(2^{|x|^k})=O(m)$), then run $N$ on $x$, in time $O(2^{|x|^k})$. So $L'\in\mathsf{NP}=\mathsf P$, decided by a deterministic $M'$ in time $m^c$. Then $L$ is decided by: on $x$, write $x\#1^{2^{|x|^k}}$ and run $M'$; time $O(2^{c(|x|^k+1)})$, so $L\in\mathsf{EXP}$. Hence $\mathsf{NEXP}\subseteq\mathsf{EXP}$, and the other inclusion is clear.
\end{solution}

\begin{solution}{ex:dag}
\idea Cycles are easy to exhibit; their absence is the complement, and $\mathsf{NL}$ is closed under complement.
\writeup A nondeterministic log-space machine can accept the graphs with a cycle: guess a vertex $v$, then a walk from $v$ of length between $1$ and $n$, vertex by vertex with a counter, and accept if it returns to $v$. A graph has a cycle if and only if some vertex returns to itself in at most $n$ steps. So the graphs with a cycle form a language in $\mathsf{NL}$, and its complement, the graphs without a cycle, is in $\mathsf{coNL}=\mathsf{NL}$ by Theorem~\ref{thm:immerman}.
\end{solution}

\begin{solution}{ex:twosatnl}
\idea Unsatisfiability is a pair of paths in the implication graph, which a log-space machine can guess.
\writeup Each clause $a\lor b$ gives the arrows $\neg a\to b$ and $\neg b\to a$ in the implication graph on the $2n$ literals. By the theorem of Aspvall, Plass and Tarjan (Tome~C1, Theorem~1.17), the formula is unsatisfiable if and only if some variable $x$ has $x$ and $\neg x$ in the same strongly connected component: $\neg x$ reachable from $x$ and $x$ from $\neg x$. A nondeterministic log-space machine guesses $x$ and the two paths, computing arrows from the clauses on demand. So unsatisfiable 2-CNF formulas form a language in $\mathsf{NL}$, and 2-SAT is in $\mathsf{coNL}=\mathsf{NL}$.
\end{solution}

\begin{solution}{ex:bipartite}
\idea Non-bipartiteness is witnessed by an odd closed walk, short enough to guess with a counter.
\writeup A graph is bipartite if and only if it has no cycle of odd length (Tome~C1, Chapter~5). If it has an odd closed walk, it has an odd cycle: split the walk at a repeated vertex into two shorter closed walks, one of which is odd, and repeat. An odd closed walk of length at most $2n+1$ exists if there is an odd cycle. A nondeterministic log-space machine guesses a vertex and a closed walk of odd length at most $2n+1$, keeping the start, the current vertex, and a counter. So non-bipartite graphs form a language in $\mathsf{NL}$, and bipartite graphs one in $\mathsf{coNL}=\mathsf{NL}$. (By Reingold's theorem it is even in $\mathsf L$.)
\end{solution}

\begin{solution}{ex:bgsconp}
\idea Diagonalise against nondeterministic machines: an accepting computation asks few questions, so freezing them leaves room for a word that changes the answer.
\writeup Let $K_C=\{1^n:\text{all words of length }n\text{ are in }C\}$. It is in $\mathsf{coNP}^C$: its complement is in $\mathsf{NP}^C$ (guess a word of length $n$ and check that it is not in $C$). We build $C$ with $K_C\notin\mathsf{NP}^C$. List the nondeterministic oracle machines $N_i$ with clocks $n^{k_i}+k_i$. At stage $i$, choose $n$ larger than all decided words with $2^n>n^{k_i}+k_i$, and consider $N_i$ on $1^n$ with the oracle that contains all undecided words of length $n$ (and none of the other undecided words). If some computation accepts, fix one; it asks fewer than $2^n$ queries. Decide the queried words according to that oracle, and put outside $C$ one word of length $n$ which it did not query; the accepting computation is still valid for the final $C$, so $N_i$ accepts $1^n$, while $1^n\notin K_C$. If no computation accepts, put all words of length $n$ into $C$; the oracle used was exactly that on length $n$, so $N_i$ rejects $1^n$ for the final $C$ as well (decide the other queried words accordingly), while $1^n\in K_C$. Either way $N_i$ errs on $1^n$, so $K_C\notin\mathsf{NP}^C$ and $\mathsf{NP}^C\ne\mathsf{coNP}^C$. (Note that in the second case all queries of all computations must be frozen; there may be exponentially many, but they all have length at most $n^{k_i}+k_i$, and the next stage chooses a larger length.)
\end{solution}

\begin{solution}{ex:pnotlinspace}
\idea Padding transports an equality of classes upwards for space but not for $\mathsf P$, which is closed under polynomial padding; the two behave differently, so they differ.
\writeup Suppose $\mathsf P=\mathsf{DSPACE}(n)$. Let $L\in\mathsf{DSPACE}(n^2)$ and $L'=\{x\#1^{|x|^2}:x\in L\}$. Then $L'\in\mathsf{DSPACE}(m)$ where $m$ is the length of the padded input: check the format and run the machine for $L$ on $x$ in space $O(|x|^2)=O(m)$. So $L'\in\mathsf P$, decided in time $m^c$. Then $L\in\mathsf P$: on $x$, write the padded word, of length $O(|x|^2)$, and run the machine for $L'$. So $\mathsf{DSPACE}(n^2)\subseteq\mathsf P=\mathsf{DSPACE}(n)$, contradicting Theorem~\ref{thm:spacehierarchy}. Hence $\mathsf P\ne\mathsf{DSPACE}(n)$.
\end{solution}

\section*{Sources and further reading}

\begin{itemize}
\item \emph{A research paper to read.} J.~Hartmanis and R.~E.~Stearns, \emph{On the computational complexity of algorithms}, Transactions of the American Mathematical Society 117 (1965), 285--306. The paper that named the field. Read Sections~1 and~2 for the definitions: they count the steps of multitape machines that \emph{compute sequences}, printing the digits of a real number one after the other, and the class $S_T$ of sequences computable in time $T(n)$ is our $\mathsf{DTIME}$ in disguise. Their Theorem~9 is the time hierarchy theorem, with the hypothesis $\inf T(n)^2/U(n)=0$, exactly the $f^2=o(g)$ of Theorem~\ref{thm:timehierarchy}, and for the same reason: their simulation of several tapes on one costs a square. Section~4 asks which real numbers are computable in real time and conjectures that algebraic irrationals are not; that question is still open.
\item \emph{The stories.} A.~Cobham, \emph{The intrinsic computational difficulty of functions}, in Proceedings of the 1964 International Congress for Logic, Methodology and Philosophy of Science, North-Holland, 1965, 24--30; J.~Edmonds, \emph{Paths, trees, and flowers}, Canadian Journal of Mathematics 17 (1965), 449--467; S.-Y.~Kuroda, \emph{Classes of languages and linear-bounded automata}, Information and Control 7 (1964), 207--223; R.~J.~Lipton, \emph{We all guessed wrong}, blog post, 2009, on the reception of Immerman and Szelepcsényi's theorem.
\item R.~E.~Stearns, J.~Hartmanis and P.~M.~Lewis~II, \emph{Hierarchies of memory limited computations}, IEEE Conference Record on Switching Circuit Theory and Logical Design (1965), 179--190; F.~C.~Hennie and R.~E.~Stearns, \emph{Two-tape simulation of multitape Turing machines}, Journal of the ACM 13 (1966), 533--546; J.~I.~Seiferas, M.~J.~Fischer and A.~R.~Meyer, \emph{Separating nondeterministic time complexity classes}, Journal of the ACM 25 (1978), 146--167.
\item W.~J.~Savitch, \emph{Relationships between nondeterministic and deterministic tape complexities}, Journal of Computer and System Sciences 4 (1970), 177--192; N.~Immerman, \emph{Nondeterministic space is closed under complementation}, SIAM Journal on Computing 17 (1988), 935--938; R.~Szelepcsényi, \emph{The method of forced enumeration for nondeterministic automata}, Acta Informatica 26 (1988), 279--284; O.~Reingold, \emph{Undirected connectivity in log-space}, Journal of the ACM 55 (2008), article 17.
\item T.~Baker, J.~Gill and R.~Solovay, \emph{Relativizations of the $\mathcal P=?\mathcal{NP}$ question}, SIAM Journal on Computing 4 (1975), 431--442; S.~Aaronson and A.~Wigderson, \emph{Algebrization: a new barrier in complexity theory}, ACM Transactions on Computation Theory 1 (2009), article 2.
\item M.~Sipser, \emph{Introduction to the Theory of Computation}, 3rd edition, Cengage, 2013, Chapters~7--9; S.~Arora and B.~Barak, \emph{Computational Complexity: A Modern Approach}, Cambridge University Press, 2009, Chapters~1--4.
\end{itemize}
